\chapter{Words, collisions, and the two units}
\label{v4-ch:words-collisions-units}

Consider three ordered positions \((x,y,x)\), with \(x\ne y\).
Neither cut \(1|23\) nor cut \(12|3\) separates the positions
on its two sides. The obstruction is the equality of the first
and last positions. A construction based on consecutive cuts
must therefore retain a nonconsecutive collision. At four
positions, the resulting support has two components, and their
coefficients meet in a nonsplit extension. These examples determine
the geometric data needed before a bar construction can be compared
with fields.

We first construct the tensor product associated with ordered cuts.
Its unit is supported at the empty configuration. The local vacuum
is a different object, present at every number of positions. We then
compute the collision supports and their coefficient complexes, and
compare a fixed word with the space of all local orders. Finally,
we separate these position spaces from the discrete set of arithmetic
places.

\section{The tensor product of ordered cuts}
\label{v4-sec:word-cut-tensor}

Let \(k\) be a field of characteristic zero and \(X\) a nonempty
smooth separated curve of finite type over \(k\). Write
\(\mathcal D(Y)\) for the unbounded derived category of right
algebraic differential modules on a smooth scheme \(Y\), with its
stable enhancement; \(\mathcal D(k)=\mathcal D(\operatorname{Spec}k)\)
is the derived category of \(k\)-vector spaces. All inverse images, direct images and tensor
products below are derived. We use the coherent differential-module
formalism of external products, extraordinary inverse images \(f^!\),
direct images \(f_*\), and localization for an open immersion.
Thus base-change isomorphisms include their compatibility with
composition. This is the formalism used in chiral Koszul duality.\footnote{
J.~Francis and D.~Gaitsgory, \emph{Chiral Koszul duality},
Selecta Mathematica \textbf{18} (2012), 27--87,
\S1.4; \href{https://arxiv.org/abs/1103.5803}{arXiv:1103.5803}.}
Our cohomological convention is
\(H^q(M[s])=H^{q+s}(M)\).

Put \([n]=\{1,\ldots,n\}\). A monotone surjection
\(\alpha:[n]\twoheadrightarrow[m]\) defines a diagonal map
\[
 d_\alpha:X^m\longrightarrow X^n,
 \qquad (y_1,\ldots,y_m)\longmapsto
                  (y_{\alpha(1)},\ldots,y_{\alpha(n)}).
\]
A \emph{word coefficient system} consists of \(F_0\in\mathcal D(k)\),
objects \(F_n\in\mathcal D(X^n)\) for \(n>0\), and coherent
equivalences \(d_\alpha^!F_n\simeq F_m\). Coherent means that
these equivalences and all their composition homotopies are part of
the datum. Its category is
\[
 \mathcal W_X=
 \mathcal D(k)\times
 \varprojlim_{\substack{[n]\ne\varnothing\\
                        \text{monotone surjections}}}
                       \mathcal D(X^n),
\]
where a surjection acts by \(d_\alpha^!\).
The empty configuration is a separate factor.

For nonnegative \(p,q\) with \(p+q=n\), let
\(U_{p,q}\subset X^n\) be the open locus
\(x_i\ne x_j\) whenever \(i\le p<j\), and write
\(j_{p,q}:U_{p,q}\hookrightarrow X^n\). If \(p=0\) or \(q=0\),
this is the identity open immersion. Define the \emph{cut product} by
\begin{equation}\label{v4-eq:word-cut-product}
 (F\star_{\mathrm{cut}}G)_n
   =\bigoplus_{p+q=n}
      j_{p,q,*}j_{p,q}^!(F_p\boxtimes G_q).
\end{equation}
Here \(X^0=\operatorname{Spec}k\), so for \(n>0\) the empty terms are the
ordinary coefficient tensors \(F_0\otimes_k G_n\) and
\(F_n\otimes_k G_0\).

\begin{proposition}[Cut product]\label{v4-prop:word-cut-product}
Formula~\eqref{v4-eq:word-cut-product} defines an associative
monoidal structure on \(\mathcal W_X\). Its unit is
\[
 K_0=k,\qquad K_n=0\quad(n>0).
\]
For \(r\ge1\) systems \(F^{(1)},\ldots,F^{(r)}\), its iterated product in
arity \(n\) is the sum over consecutive ordered decompositions of
\([n]\) into \(r\) blocks, localized where positions in different
blocks are distinct.
\end{proposition}

\begin{proof}
Let \(j:U\hookrightarrow Y\) be an open immersion with closed
complement \(i:Z\hookrightarrow Y\). Localization gives
\(i_*i^!M\to M\to j_*j^!M\), and \(i^!j_*=0\).
For \(N\in\mathcal D(U)\) and \(f:Y'\to Y\), transitivity of
extraordinary inverse images
therefore shows that \(f^!j_*N\) has zero extraordinary restriction
to \(f^{-1}Z\). Its restriction to \(f^{-1}U\) is the pullback
of \(N\). Localization and the full faithfulness of extension
from an open set identify it with
\(j'_*f_U^!N\), where \(j':f^{-1}U\hookrightarrow Y'\).
These identifications are coherent under composition.

Apply this to \(f=d_\alpha\) and \(U=U_{p,q}\). If the cut
separates two labels in the same fibre of \(\alpha\), its inverse
image is empty. Otherwise it is the inverse image of a unique cut
of \([m]\), and the coefficient equivalences for \(F\) and \(G\)
give the corresponding summand in arity \(m\). Thus the displayed
sum is a word coefficient system.

In an iterated product, successive localizations require exactly
that positions in distinct final blocks be disjoint. Each
parenthesization therefore gives the same sum, with the same
coefficient and open immersion in every summand. The coherent
base-change maps identify these sums; their further compositions
agree on the same final open set. This proves every associativity
coherence. If one input is \(K\), only its empty block survives,
and the remaining summand is the other input. These identifications
give the unit coherences.
\end{proof}

Call \(F\) \emph{positive} if \(F_0=0\). For positive inputs,
all blocks in the proposition are nonempty. There are
\(\binom{n-1}{r-1}\) such decompositions when \(1\le r\le n\),
and none when \(r>n\): choose their \(r-1\) cuts among the
\(n-1\) gaps. This finiteness concerns the number of geometric
blocks. It places no bound on an auxiliary bar length in a
coefficient algebra at the empty configuration.

\section{The empty unit and the local vacuum}
\label{v4-sec:word-two-units}

There is also a pointwise tensor product of differential modules,
using the diagonal \(\Delta_{X^n}:X^n\hookrightarrow X^n\times X^n\):
\[
 (F\otimes^!G)_n=\Delta_{X^n}^!(F_n\boxtimes G_n).
\]
Its unit is the dualizing system
\begin{equation}\label{v4-eq:word-vacuum}
 J_0=k,\qquad J_n=p_n^!k=\omega_{X^n}[n]\quad(n>0),
\end{equation}
where \(p_n:X^n\to\operatorname{Spec}k\).
Transitivity gives \(d_\alpha^!J_n=J_m\). The shift in
\eqref{v4-eq:word-vacuum} is part of the right differential-module
convention; \(J_n\) is not the unshifted canonical bundle.

A counital coassociative cut coalgebra is an object with a
coproduct \(C\to C\star_{\mathrm{cut}}C\) and a counit
\(C\to K\), satisfying their coherent identities. A coaugmentation
is a coalgebra map \(K\to C\). We call it connected when
\(C_0=k\) and the empty components of counit and coaugmentation
are identities.

\begin{proposition}[Vacuum coalgebra]\label{v4-prop:word-vacuum-coalgebra}
The system \(J\) is a connected cut coalgebra. Its component at a
cut is the localization map
\begin{equation}\label{v4-eq:word-vacuum-coproduct}
 J_n\longrightarrow j_{p,q,*}J_{U_{p,q}}
       =j_{p,q,*}j_{p,q}^!(J_p\boxtimes J_q).
\end{equation}
The maps \(K\to J\to K\) are identity at the empty
configuration and zero in positive arity.
\end{proposition}

\begin{proof}
The adjunction for \(j_{p,q}\) gives the displayed map.
Two successive coproducts localize on the locus where three
consecutive blocks are pairwise disjoint. Both orders of
localization give its adjunction map. The same argument works for
every number of blocks, proving the coherent coalgebra identities.
The cuts with an empty block are identities. They prove the
counit and coaugmentation identities.
\end{proof}

A \emph{local vacuum marking} of a connected cut coalgebra is a
specified coalgebra map \(u:J\to C\). The coaugmentation
\(K\to C\) determines only its empty value. Even if a nonzero
unary vacuum is specified, extending it to such a marking can fail.

\begin{example}[A unary vacuum that does not extend]
\label{v4-ex:word-vacuum-obstruction}
Let \(\Delta_n:X\hookrightarrow X^n\) be the small diagonal and put
\[
 P_0=k,\qquad P_n=\Delta_{n,*}J_1\quad(n>0).
\]
Closed base change and Kashiwara equivalence give
\(d_\alpha^!P_n=P_m\). Every positive cut has empty intersection
with the small diagonal. Give \(P\) zero coproduct at these cuts
and identity coproduct at each cut with an empty block. This is a
connected cut coalgebra: an iterated coproduct with two or more
nonempty output blocks is zero in every parenthesization.

There is no coalgebra map \(u:J\to P\) with \(u_1=\mathrm{id}_{J_1}\).
Indeed, restrict its binary coproduct identity to
\(U=X^2\setminus\Delta_2\). The path through \(P_2\) is zero,
whereas the path through \(J_1\boxtimes J_1\) is the identity of
\(J_U\). This object is nonzero, since a smooth nonempty curve has
two distinct geometric points. The contradiction concerns the
specified unary identity. The map \(J\to K\to P\), whose
positive components vanish, still exists.
\end{example}

\section{Collisions that block every cut}
\label{v4-sec:word-collision-supports}

Write \(\Delta_{ij}\subset X^n\) for the equality \(x_i=x_j\).
For \(n\ge2\), define the reduced closed subscheme
\begin{equation}\label{v4-eq:word-collision-support}
 Z_n=
 \left(\bigcap_{p=1}^{n-1}
               \bigcup_{i\le p<j}\Delta_{ij}\right)_{\mathrm{red}}.
\end{equation}
Thus a tuple lies in \(Z_n\) exactly when every consecutive cut
is crossed by an equality. Put \(Z_1=X\). The first cases are
\begin{equation}\label{v4-eq:word-first-supports}
 Z_2=\Delta_{12},\qquad Z_3=\Delta_{13},\qquad
 Z_4=\Delta_{14}\cup(\Delta_{13}\cap\Delta_{24}).
\end{equation}
For the last equality, a tuple with \(x_1\ne x_4\) must still
cross its first and last gaps. If it does not have both
\(x_1=x_3\) and \(x_2=x_4\), these two requirements force
\(x_1=x_2\) and \(x_3=x_4\). The middle cut then forces all
four coordinates equal, contrary to \(x_1\ne x_4\).
The converse inclusions follow by checking the three gaps.

There is an equally explicit description in every arity. A matching
on \([n]\) is a set of disjoint pairs \((a,b)\), with \(a<b\).
Associate to it the closed subscheme
\(\Delta_M=\bigcap_{(a,b)\in M}\Delta_{ab}\).
The pair \((a,b)\) crosses the gaps \(a,a+1,\ldots,b-1\).
Call \(M\) spanning if its pairs cross every gap, and irredundant
if deleting any pair destroys this property.

\begin{theorem}[Components of the cut support]
\label{v4-thm:word-support-components}
Assume \(X\) is geometrically integral and \(n\ge2\). The irreducible components
of \(Z_n\) are the \(\Delta_M\) for irredundant spanning
matchings \(M\). A matching with \(r\) pairs has codimension \(r\),
and there are
\[
 \binom{n-2}{2r-2},\qquad 1\le r\le\lfloor n/2\rfloor,
\]
such components. For \(n\ge3\) their total number is \(2^{n-3}\).
\end{theorem}

\begin{proof}
The equalities of a geometric tuple partition \([n]\) into blocks.
In a block with at least two labels, replace all its equalities by
the equality of its smallest and largest labels. This retains
exactly the gaps crossed by that block. The resulting pairs are
disjoint. If the tuple lies in \(Z_n\), they form a spanning
matching; remove redundant pairs to obtain \(M\). Consequently
\(Z_n\) is the union of the \(\Delta_M\) in the statement.

Order the pairs by their left endpoints,
\(a_1<\cdots<a_r\). An interval contained in another is
redundant, so their right endpoints also increase. Spanning gives
\(a_1=1\), \(b_r=n\), and \(a_{i+1}<b_i\): equality of
these endpoints is excluded because this is a matching. A pair
has a gap not crossed by its neighbours exactly when
\(b_{i-1}<a_{i+1}\), for \(1<i<r\). Thus, for \(r>1\),
the endpoints occur in precisely the order
\[
 1=a_1<a_2<b_1<a_3<b_2<\cdots<a_r<b_{r-1}<b_r=n.
\]
Conversely this order makes the matching spanning and gives every
pair a private gap. Choosing its \(2r-2\) interior endpoints
therefore gives the binomial coefficient above. For \(r=1\),
the unique pair is \((1,n)\).

Each \(\Delta_M\) is isomorphic to \(X^{n-r}\), hence is
irreducible of codimension \(r\). Containment between two such
diagonals would make the partition into pairs for one refine that
for the other. Since a pair has no nontrivial proper refinement
except deletion, irredundancy excludes a strict containment.
They are therefore exactly the components. Summing the even
binomial coefficients of order \(n-2\) gives \(2^{n-3}\)
when \(n\ge3\).
\end{proof}

\section{Coefficient complexes on the collision support}
\label{v4-sec:word-support-coefficients}

A support records where a coefficient can live. It does not record
the coefficient or its maps along intersections. We make those maps
explicit for \(X=\mathbb A^1_k\). Let
\(R=k[z_1,\ldots,z_n]\). For this calculation we use complexes
of \(R\)-modules. Their localizations have natural left
differential-module actions. Tensoring with \(\omega_{X^n}[n]\)
gives the corresponding supported right-module coefficients of
\(J_n\); the scalar complexes below carry no such shift.
For \(f\in R\), put
\[
 L_f=[\,R\longrightarrow R[f^{-1}]\,],
\]
in cohomological degrees \(0,1\), with the localization differential.
The augmentation \(\epsilon_f:L_f\to R\) is identity in degree
zero and zero in degree one. This flat complex represents
\(R\Gamma_{V(f)}R\), the derived functor of sections supported
on \(V(f)\). For a general coefficient complex \(Q\),
\(L_f\otimes_R Q\) represents \(R\Gamma_{V(f)}Q\).

Define
\begin{equation}\label{v4-eq:word-cut-cech}
 g_p=\prod_{i\le p<j}(z_i-z_j),\qquad
 C_n=\bigotimes_{p=1}^{n-1}L_{g_p}.
\end{equation}
This is a bounded complex of flat modules, with \(2^{n-1}\)
localization summands. Iterated local cohomology gives
\(C_n\simeq R\Gamma_{Z_n}R\). Local cohomology depends on the
reduced support. The ideal generated by the \(g_p\) need not be
radical: at \(n=3\), with \(x=z_1-z_2\) and \(y=z_2-z_3\), it is
\[
 (x(x+y),y(x+y))=(x+y)(x,y),
\]
whose radical is \((x+y)\).

The component description gives another model that displays the
attachments. Choose a total order on the finite set \(\mathcal M_n\)
of irredundant spanning matchings. For \(M\in\mathcal M_n\), let
\[
 K_M=\bigotimes_{(a,b)\in M}L_{z_a-z_b},
\]
with pairs ordered by their left endpoints. For a nonempty subset
\(A=\{M_0<\cdots<M_p\}\), put
\(K_A=K_{M_0}\otimes_R\cdots\otimes_RK_{M_p}\).
Repeated equations remain repeated tensor factors in this formula.

\begin{proposition}[Component model]
\label{v4-prop:word-component-complex}
The complex
\begin{equation}\label{v4-eq:word-component-complex}
 S_n=\bigoplus_{\varnothing\ne A\subseteq\mathcal M_n}
                           K_A[|A|-1]
\end{equation}
represents \(R\Gamma_{Z_n}R\). On \(K_A[p]\), its internal
differential is \((-1)^p d_{K_A}\); its other differential is
\(\sum_{j=0}^p(-1)^j\partial_j\), where \(\partial_j\)
applies the augmentation of \(K_{M_j}\) and deletes that factor.
The latter differential is zero when \(p=0\). The augmentations
give a natural span of quasi-isomorphisms
\begin{equation}\label{v4-eq:word-support-roof}
 C_n\longleftarrow C_n\otimes_R S_n\longrightarrow S_n.
\end{equation}
The same assertions hold after tensoring with any complex \(Q\).
\end{proposition}

\begin{proof}
The tensor of the \(L_f\)'s computes local cohomology supported
on the intersection of their zero sets. Thus \(K_A\) represents
\(R\Gamma_{\cap_{M\in A}\Delta_M}R\). The augmentation of a
factor induces the inclusion of sections supported on this
intersection into sections supported on the intersection with that
factor omitted. These are the maps \(\partial_j\).

Two deletions commute before their alternating signs, so the
horizontal differential squares to zero. It commutes with each
unshifted internal differential; the factor \((-1)^p\) makes the
two total differentials anticommute. For two closed sets, the
resulting complex is the local-cohomology Mayer--Vietoris triangle
\[
 R\Gamma_{Z\cap Z'}R\longrightarrow
 R\Gamma_ZR\oplus R\Gamma_{Z'}R
                 \longrightarrow R\Gamma_{Z\cup Z'}R.
\]
Iterate this triangle over the finite component set. At each step
the new intersections are exactly the summands containing the new
component, with the displayed deletion maps. This proves
\eqref{v4-eq:word-component-complex}.

Local cohomology supported on \(Z_n\) is idempotent. Indeed, its
localization triangle becomes the identity triangle on an object
whose restriction to the complement of \(Z_n\) vanishes.
Apply this to \(C_n\) and \(S_n\). The two counits are exactly
the arrows of \eqref{v4-eq:word-support-roof}, so both are
quasi-isomorphisms. Every term used is flat and the complexes are
bounded. Tensoring the construction with an arbitrary complex
therefore computes its derived tensor and preserves these
quasi-isomorphisms.
\end{proof}

At the generic point of a component with \(r\) pairs, invert the
equations that exclude the other components and choose independent
normal coordinates \(u_1,\ldots,u_r\) for the pairs. Its local
cohomology is concentrated in degree \(r\), represented by
negative powers of all the \(u_i\). In particular it contains
\(1/(u_1\cdots u_r)\); the iterated residue of
\(du_1\cdots du_r/(u_1\cdots u_r)\), in that order, is one.
This follows directly by tensoring the \(r\) two-term localization
complexes, or by their Laurent monomial basis.

The idempotence just used is a property of local cohomology.
It must be distinguished from derived intersection. For example,
over \(k[u]\), the free resolution
\([\,k[u]\xrightarrow{u}k[u]\,]\) of \(k\), in degrees
\(-1,0\), gives
\begin{equation}\label{v4-eq:word-derived-self-intersection}
 k\otimes_{k[u]}^{\mathbf L}k\simeq k\oplus k[1],
 \qquad
 L_u\otimes_{k[u]}L_u\simeq L_u.
\end{equation}
The first complex retains a normal Tor class. The second retains
all principal parts along the same closed support.

\section{The four-point attachment and normal derivatives}
\label{v4-sec:word-four-point-attachment}

At four positions choose
\(a=z_1-z_4\), \(b=z_1-z_3\), \(c=z_2-z_4\), and the
spectator coordinate \(z_4\). Suppress the spectator polynomial
factor and put \(R=k[a,b,c]\). Formula~\eqref{v4-eq:word-first-supports}
becomes
\[
 Z_4=V(a)\cup V(b,c)=V(ab,ac).
\]
Its coefficient complex is
\begin{equation}\label{v4-eq:word-four-point-cech}
 R\xrightarrow{f\mapsto(f,f)}
 R[(ab)^{-1}]\oplus R[(ac)^{-1}]
 \xrightarrow{(f,g)\mapsto g-f}R[(abc)^{-1}].
\end{equation}
Its terms have cohomological degrees \(0,1,2\).

\begin{proposition}[Nonsplit attachment]
\label{v4-prop:word-four-point-nonsplit}
Let
\[
 M=R[b^{-1},c^{-1}]\big/
                  \bigl(R[b^{-1}]+R[c^{-1}]\bigr).
\]
The only nonzero cohomology of
\eqref{v4-eq:word-four-point-cech} is
\[
 H^1_{Z_4}(R)=R[a^{-1}]/R,\qquad
 H^2_{Z_4}(R)=M[a^{-1}].
\]
The attachment of \(V(b,c)\) to the union is the nonsplit exact
sequence of \(R\)-modules
\begin{equation}\label{v4-eq:word-four-point-extension}
 0\longrightarrow M\longrightarrow M[a^{-1}]
   \longrightarrow H^3_{(a,b,c)}(R)\longrightarrow0.
\end{equation}
\end{proposition}

\begin{proof}
Inside \(R[(abc)^{-1}]\), the two middle localizations intersect
in \(R[a^{-1}]\). This computes \(H^1\). Their quotient in
the last term has basis \(a^i b^{-j}c^{-l}\), with
\(i\in\mathbb Z\) and \(j,l>0\), which gives \(M[a^{-1}]\).
There is no degree-zero kernel because localization of the
polynomial ring is injective.

The module \(M\) is free over \(k[a]\), with basis
\(b^{-j}c^{-l}\) for \(j,l>0\). Its quotient after inverting
\(a\) has basis with all three exponents negative. This is the
top local cohomology at \((a,b,c)\), proving exactness. Every
element of this quotient is killed by a power of \(a\), whereas
multiplication by \(a\) is invertible on \(M[a^{-1}]\).
Every \(R\)-linear map from the quotient back to the middle term
is consequently zero. The quotient contains the nonzero class
\(1/(abc)\), so the surjection has no section.
\end{proof}

The intersection \(V(a,b,c)\) cannot therefore be replaced by an
independent direct summand while preserving the coefficient action.
Even at one normal coordinate, a scalar residue loses this action.
To see it in the right differential-module convention, set
\[
 P=k[u,u^{-1}]\,du\big/k[u]\,du,
 \qquad e_r=[u^{-r-1}du]\quad(r\ge0).
\]
The right Weyl action is
\begin{equation}\label{v4-eq:word-normal-derivatives}
 e_0u=0,\qquad e_ru=e_{r-1}\ (r>0),\qquad
 e_r\partial_u=(r+1)e_{r+1}.
\end{equation}
Indeed, right differentiation of a density \(f\,du\) is
\(-f'\,du\). The formulas satisfy
\(\partial_u u-u\partial_u=1\) and give
\(e_0\partial_u^r=r!e_r\). All normal derivatives are present.
The scalar residue sends \(e_0\) to one and \(e_r\) to zero
for \(r>0\), but its kernel is not stable under \(u\):
\(\operatorname{Res}(e_1u)=1\). It is therefore insufficient
as a replacement for this coefficient module.

\section{The outer residue}
\label{v4-sec:word-outer-residue}

The nonconsecutive support already changes a three-position
calculation. Put \(x=z_1-z_2\), \(y=z_2-z_3\), and consider
\[
 f(x,y)=\frac{1}{x(x+y)}.
\]
Let \(L\) be the coefficient of \(x^{-1}y^{-1}\) in the
expansion for \(|x|<|y|\), and \(R\) the same iterated
coefficient for \(|y|<|x|\). These inequalities specify formal
Laurent expansions; no analytic limiting operation is required.
For the outer collision use \(u=x+y=z_1-z_3\). Let \(N\)
be the coefficient of \(u^{-1}y^{-1}\) in the expansion
for \(|u|<|y|\). The volume form is unchanged:
\(dx\wedge dy=du\wedge dy\).

\begin{proposition}\label{v4-prop:word-outer-residue}
With these coordinates and orientations,
\[
 L=1,\qquad R=0,\qquad N=-1,\qquad L-R+N=0.
\]
\end{proposition}

\begin{proof}
The three expansions are respectively
\[
 \sum_{r\ge0}(-1)^r x^{r-1}y^{-r-1},\qquad
 \sum_{r\ge0}(-1)^r y^r x^{-r-2},\qquad
 -\sum_{r\ge0}u^{r-1}y^{-r-1}.
\]
Their indicated coefficients are \(1,0,-1\).
\end{proof}

This residue identity does not by itself construct an associativity
homotopy. In a specified operation complex \(Q\), closed
degree-zero operations \(L,R,N\) satisfying \(L-R+N=0\)
admit \(h\in Q^{-1}\) with \(dh=L-R=-N\) exactly when
\([N]=0\) in \(H^0(Q)\). For \(Q=k\) in degree zero,
the three scalars just computed satisfy the identity and admit
no such \(h\). A field comparison must specify the operation
complex, its supported coefficients and the primitive there.

\section{All local orders and the word comparison}
\label{v4-sec:word-ordered-atlas}

A fixed word has consecutive cuts. To retain every local ordering
at a collision, let \(I\) be a nonempty finite set and take a
copy \(X^I_\sigma\) of \(X^I\) for each total order \(\sigma\)
on \(I\). For two orders define
\[
 V_{\sigma,\rho}=
 X^I\setminus\bigcup_{\substack{i,j\in I\\
              \sigma\text{ and }\rho\text{ order }i,j\text{ differently}}}
                          \Delta_{ij}.
\]
Glue the two copies by the identity on this open subset, and denote
the resulting scheme by \(\mathscr R_I\).
The chart coordinates induce \(q_I:\mathscr R_I\to X^I\).
For \(X=\mathbb A^1_{\mathbb C}\), this is the configuration
space used for raviolo vertex algebras.\footnote{
L.~Alfonsi, H.~Kim and C.~A.~S.~Young,
\emph{Raviolo vertex algebras, cochains and conformal blocks},
\S4.1; \href{https://arxiv.org/abs/2401.11917}{arXiv:2401.11917}.}

\begin{proposition}[Ordered configuration atlas]
\label{v4-prop:word-ordered-atlas}
These identifications satisfy the gluing cocycle. The map
\(q_I:\mathscr R_I\to X^I\) is a local isomorphism, and
\(\mathscr R_I\) is smooth and quasi-compact with quasi-compact
intersections of affine opens. Over a geometric tuple whose equality
blocks are \(B\), its fibre has
\[
                    \prod_B |B|!
\]
points. In particular the fibre at a threefold collision has six
points.
\end{proposition}

\begin{proof}
On a triple overlap, if two orders disagree on a pair, at least
one of their comparisons with the third order disagrees on that
pair. Thus
\(V_{\sigma,\rho}\cap V_{\rho,\tau}\subseteq V_{\sigma,\tau}\).
All transition maps are identities, which proves the cocycle.
The charts give smoothness and the local-isomorphism assertion.
There are finitely many charts; their intersections are complements
of finitely many diagonals in the noetherian scheme \(X^I\),
so the asserted quasi-compactness follows.

Over a fixed tuple, two charts are identified exactly when they
give the same order within each equality block. Every choice of
orders on these blocks extends to a total order on \(I\).
The choices on different blocks are independent, proving the
fibre formula.
\end{proof}

A surjection \(\alpha:I\twoheadrightarrow J\) with an order
on each fibre defines \(\delta_\alpha:\mathscr R_J\to\mathscr R_I\).
On a chart ordered by \(\tau\) on \(J\), expand each label
to its ordered fibre and use the usual diagonal map of coordinates.
The inverse image of an overlap is precisely the overlap for the
two orders on \(J\), so these chart maps glue. Composing such
surjections concatenates the ordered fibres and gives the same
diagonal map. These collision maps need not be closed immersions.
For the ordered fibres \((1,3)|(2)\), their image in the chart
\(1<2<3\) is \(\Delta_{13}\setminus\Delta_{123}\), which
is not closed in that chart.

Define \(\mathcal C_X\) to be the category of coherent systems
\(F_I\in\mathcal D(\mathscr R_I)\) with
\(\delta_\alpha^!F_I\simeq F_J\), together with a separate
empty coefficient. Differential modules on this possibly
nonseparated scheme are defined by descent from its charts.
For an ordered decomposition \(I=I_1\sqcup I_2\), let \(U_{I_1,I_2}\)
be the locus of disjoint positions in different blocks. On that
locus, orders between the blocks are identified by the gluing,
so the atlas is the corresponding open part of
\(\mathscr R_{I_1}\times\mathscr R_{I_2}\). Define \(F\star G\)
by summing the localized external products over all such ordered
decompositions, including empty blocks. The proof of the cut-product
proposition applies: a decomposition splitting a collision fibre
has empty pullback; every surviving decomposition descends to the
quotient. All iterates localize on the same disjoint-block locus.
This gives a symmetric monoidal product, whose symmetry exchanges
the two blocks.

Restriction to the standard chart \(1<\cdots<n\) defines
\(\operatorname{res}:\mathcal C_X\to\mathcal W_X\).

\begin{proposition}[Projection to consecutive cuts]
\label{v4-prop:word-atlas-comparison}
There are natural maps
\begin{align*}
 \iota_{F,G}:\operatorname{res}F\star_{\mathrm{cut}}
                    \operatorname{res}G
       &\longrightarrow\operatorname{res}(F\star G),\\
 \pi_{F,G}:\operatorname{res}(F\star G)
       &\longrightarrow\operatorname{res}F\star_{\mathrm{cut}}
                    \operatorname{res}G,
 \qquad \pi_{F,G}\iota_{F,G}=\mathrm{id}.
\end{align*}
The maps \(\iota\) make restriction lax monoidal, and the maps
\(\pi\) make it colax monoidal. They transport algebra and
coalgebra structures, respectively. Neither comparison is an
equivalence in general.
\end{proposition}

\begin{proof}
On a standard chart the atlas product is the finite direct sum
over all ordered decompositions of \([n]\). Include the summands
whose first block is a prefix, or project onto those summands.
Under a monotone collision, a surviving prefix is exactly the
inverse image of a quotient prefix. The maps therefore respect
the word coefficient equivalences.

Every iterated inclusion or projection selects the same final
consecutive decomposition with the same localized coefficient.
This proves their monoidal coherences, including those for empty
blocks, and projection after inclusion is identity. For two
positive systems equal to the dualizing system in positive arity,
the separated two-point fibre of the atlas product is
\(J_U\oplus J_U\), indexed by \(1|2\) and \(2|1\).
The cut product has only the \(1|2\) summand. Thus \(\pi\)
has a nonzero kernel there, and \(\iota\) has a nonzero
cokernel.
\end{proof}

With \(n\) positive inputs at \(n\) distinct positions, the
same calculation gives \(n!\) atlas summands and one cut
summand. A bar--cobar theorem for the cut product therefore does
not identify the corresponding construction on the full atlas.
Any such comparison must retain its actual maps on the omitted
orders and their collision coefficients.

\section{Arithmetic places and position geometry}
\label{v4-sec:word-arithmetic-positions}

An arithmetic place is a label for a local coefficient object.
It need not be a variable position on the curve carrying a field.
This distinction already appears in the simplest functors of points.
For every scheme \(S\), there is exactly one map
\(S\to\operatorname{Spec}\mathbb Z\). Consequently the functor
of nonempty finite subsets of
\(\operatorname{Spec}\mathbb Z(S)\) is terminal. The closed
points \(\operatorname{Spec}\mathbb F_p\to\operatorname{Spec}\mathbb Z\)
use different test schemes; they are not several elements of this
set for one fixed \(S\).

Let \(\mathbb P\) be the set of rational primes and
\(V=\mathbb P\sqcup\{\infty\}\) the abstract set of places.
Its nonempty finite subsets and its nonempty finite multisets are
the discrete sets
\[
 \operatorname{Fin}_{+}(V),\qquad
 \operatorname{Mult}_{+}(V)=\coprod_{n\ge1}V^n/\mathfrak S_n.
\]
The coproduct in the second formula records multiplicity; it
does not require transition maps between different symmetric powers.
For a commutative ring \(A\), let \(\mathcal D(A)\) here denote
the derived category of \(A\)-modules. Given complexes
\(H_v\in\mathcal D(A)\), finite tensor products give
\begin{equation}\label{v4-eq:word-discrete-factorization}
 F(S)=\bigotimes_{v\in S,A}^{\mathbf L}H_v,\qquad
 F(S\sqcup T)\simeq F(S)\otimes_A^{\mathbf L}F(T)
                         \quad(S\cap T=\varnothing).
\end{equation}
We include the empty set by \(F(\varnothing)=A\).
The associators and symmetries of derived tensor give all
coherences, with the usual Koszul signs on complexes. This supplies
factorization for disjoint labels. Maps at coincident positions,
normal derivatives and residue operations require a specified
position geometry and further coefficient maps.

To assemble infinitely many places, units must also be specified.
More generally, let \(\mathcal E\) be a cocomplete symmetric monoidal
\(\infty\)-category with unit \(\mathbf1\), and let \(H_v\in\mathcal E\).
Choose a cofinite subset \(V^{\mathrm{sph}}\subseteq V\) and maps
\(u_v:\mathbf1\to H_v\) for \(v\in V^{\mathrm{sph}}\).
Let \(\mathrm{Fin}^{\mathrm{sph}}(V)\) be the poset of finite
subsets containing \(V\setminus V^{\mathrm{sph}}\), ordered by inclusion.
For \(S\subseteq T\), tensor with the maps \(u_v\) at the new
places to obtain
\[
 \iota_{S,T}:\bigotimes_{v\in S}H_v
           \longrightarrow\bigotimes_{v\in T}H_v.
\]
The colimit
\begin{equation}\label{v4-eq:word-restricted-tensor}
 \bigotimes_{v\in V}'(H_v,u_v)
   =\varinjlim_{S\in\mathrm{Fin}^{\mathrm{sph}}(V)}
                                      \bigotimes_{v\in S}H_v
\end{equation}
is called the \emph{Tate restricted tensor product} for these
specified reference maps. It is a colimit of a covariant diagram.

Indeed, \(\iota_{T,U}\iota_{S,T}=\iota_{S,U}\) coherently:
both maps insert the same \(u_v\)'s in the same labelled factors,
and tensor coherence identifies their parenthesizations. The
index poset is nonempty and filtered, since the exceptional set
is finite and the union of two allowed finite sets is allowed.
This constructs the diagram and its colimit. To obtain an algebra,
coalgebra or monoidal category by this construction, its reference
maps must preserve that structure and its colimit must be taken
in a coefficient category that supports the required operations.

\begin{proposition}[Three finite-place coefficient models]
\label{v4-prop:word-finite-place-models}
Let \(p\) be a rational prime. Define the completed tensor below
to mean the ordinary \(p\)-adic completion. Then
\begin{align}
 \varprojlim_r
  (\mathbb Z_p\otimes_{\mathbb Z}\mathbb Z_p)/p^r
       &\simeq\mathbb Z_p,\label{v4-eq:word-padic-tensor}\\
 \mathbb F_p\otimes_{\mathbb Z}^{\mathbf L}\mathbb F_p
       &\simeq\mathbb F_p\oplus\mathbb F_p[1]
          \quad\text{as a complex},\label{v4-eq:word-prime-tor}\\
 \mathbb Z_p[[t_1,t_2]]/(t_1-t_2)
       &\simeq\mathbb Z_p[[t]].\label{v4-eq:word-padic-position}
\end{align}
The kernel in the last line is the nonzero principal ideal
\((t_1-t_2)\).
\end{proposition}

\begin{proof}
Modulo \(p^r\), the first tensor product is
\((\mathbb Z/p^r)\otimes_{\mathbb Z}\mathbb Z_p
 =\mathbb Z/p^r\); multiplication induces this isomorphism
compatibly for all \(r\). Taking the inverse limit proves
\eqref{v4-eq:word-padic-tensor}. The free resolution
\([\,\mathbb Z\xrightarrow{p}\mathbb Z\,]\) of \(\mathbb F_p\),
in degrees \(-1,0\), becomes a complex with zero differential
after tensoring with \(\mathbb F_p\). This proves
\eqref{v4-eq:word-prime-tor}. Finally the formal change of
variables \(u=t_1-t_2\), \(t=t_2\) identifies the last ring
before quotienting with \(\mathbb Z_p[[u,t]]\). Evaluation at
\(u=0\) has kernel \((u)\), proving the last assertion.
\end{proof}

The completed base tensor, the derived closed-point intersection
and the formal position disc thus retain different data. The
first identity does not eliminate the normal class in the second
or the position variable in the third. A finite-place field
construction can use a chosen formal disc; crystalline or
prismatic structures are further possible choices when their
Frobenius data are required. None is forced solely by the discrete
place index or by the first tensor identity. The arithmetic
comparison must specify which position space and which coefficient
operations its maps preserve.
